%% Đề 47 -- de-tuyen-sinh-lop-10-thpt-mon-toan-nam-2021-2022-so-gddt-lang-son.pdf
\deheader{ĐỀ SỐ 47}%
  {Kỳ thi Tuyển sinh vào lớp 10 THPT — Năm học 2021--2022}%
  {Sở GD\&ĐT Lạng Sơn | Chỉ phần Đại số}
\addcontentsline{toc}{section}{Đề 47 — Tuyển sinh vào lớp 10 THPT — Năm học 2021--2022}

\begin{de_cau}{Câu 1}{2,0 điểm}
\begin{enumerate}[label=\alph*)]
  \item Tính: $A = \sqrt{(2-\sqrt{3})^2}+\sqrt{3}$.
  \item Rút gọn: $B = \left(\dfrac{\sqrt{x}}{\sqrt{x}-1}+\dfrac{1}{\sqrt{x}+1}\right):\dfrac{\sqrt{x}+1}{\sqrt{x}}$ với $x>0$, $x\neq1$.
\end{enumerate}
\end{de_cau}

\begin{loi_giai_cau}
\textbf{a)} $\sqrt{(2-\sqrt{3})^2}=|2-\sqrt{3}|=2-\sqrt{3}$ (vì $2>\sqrt{3}$).

$A=(2-\sqrt{3})+\sqrt{3}=\boxed{2}$.

\textbf{b)} $\dfrac{\sqrt{x}(\sqrt{x}+1)+(\sqrt{x}-1)}{(\sqrt{x}-1)(\sqrt{x}+1)}=\dfrac{x+\sqrt{x}+\sqrt{x}-1}{x-1}=\dfrac{x+2\sqrt{x}-1}{x-1}$.

Nhân với $\dfrac{\sqrt{x}}{\sqrt{x}+1}$:

$B=\dfrac{(x+2\sqrt{x}-1)\sqrt{x}}{(x-1)(\sqrt{x}+1)}=\dfrac{\sqrt{x}(x+2\sqrt{x}-1)}{(\sqrt{x}-1)(\sqrt{x}+1)^2}$.

Cách đơn giản hơn:
$\dfrac{\sqrt{x}}{\sqrt{x}-1}+\dfrac{1}{\sqrt{x}+1}=\dfrac{x+\sqrt{x}+\sqrt{x}-1}{x-1}=\dfrac{x+2\sqrt{x}-1}{x-1}$.

$B=\dfrac{x+2\sqrt{x}-1}{x-1}\cdot\dfrac{\sqrt{x}}{\sqrt{x}+1}=\dfrac{\sqrt{x}(x+2\sqrt{x}-1)}{(\sqrt{x}+1)^2(\sqrt{x}-1)}=\boxed{\dfrac{\sqrt{x}(x+2\sqrt{x}-1)}{(x-1)(\sqrt{x}+1)}}$.
\end{loi_giai_cau}

\begin{de_cau}{Câu 2}{2,0 điểm}
\begin{enumerate}[label=\alph*)]
  \item Giải phương trình: $x^2-4x+3=0$.
  \item Giải hệ phương trình:
  \[\begin{cases} x+2y=4 \\ 3x-y=1 \end{cases}.\]
\end{enumerate}
\end{de_cau}

\begin{loi_giai_cau}
\textbf{a)} $(x-1)(x-3)=0\Rightarrow\boxed{x=1}$ hoặc $\boxed{x=3}$.

\textbf{b)} Từ (1): $x=4-2y$. Thế vào (2): $3(4-2y)-y=1\Rightarrow12-7y=1\Rightarrow y=\dfrac{11}{7}$, $x=4-\dfrac{22}{7}=\dfrac{6}{7}$.

Nghiệm $\boxed{(x;y)=\left(\dfrac{6}{7};\dfrac{11}{7}\right)}$.
\end{loi_giai_cau}

\begin{de_cau}{Câu 3}{1,5 điểm}
Một người đi bộ từ A đến B. Nếu đi với vận tốc 5 km/h thì đến sớm hơn dự định 12 phút, nếu đi với vận tốc 4 km/h thì đến muộn hơn dự định 18 phút. Tính độ dài quãng đường AB.
\end{de_cau}

\begin{loi_giai_cau}
Gọi $d$ = độ dài AB (km), $t$ = thời gian dự định (giờ).

\[\frac{d}{5}=t-\frac{1}{5};\qquad\frac{d}{4}=t+\frac{3}{10}.\]
Trừ vế: $\dfrac{d}{4}-\dfrac{d}{5}=\dfrac{1}{2}\Rightarrow\dfrac{d}{20}=\dfrac{1}{2}\Rightarrow d=\boxed{10}$ km.
\end{loi_giai_cau}

\begin{de_cau}{Câu 4}{0,5 điểm}
Cho phương trình $x^2-2mx+m^2-m=0$ (tham số $m$). Tìm $m$ để phương trình có hai nghiệm dương phân biệt $x_1, x_2$ thỏa mãn $x_1^2+x_2^2=9$.
\end{de_cau}

\begin{loi_giai_cau}
$\Delta'=m^2-(m^2-m)=m>0$.

Vi-ét: $x_1+x_2=2m>0$ (\checkmark ), $x_1x_2=m^2-m=m(m-1)>0\Rightarrow m>1$.

$x_1^2+x_2^2=4m^2-2m(m-1)=2m^2+2m=9\Rightarrow2m^2+2m-9=0\Rightarrow m=\dfrac{-2+\sqrt{22}}{2}=-1+\dfrac{\sqrt{22}}{2}$.

Kiểm tra $m>1$: $\dfrac{\sqrt{22}}{2}\approx2.35>2$, nên $m\approx1.35>1$ \checkmark .

Vậy $\boxed{m=\dfrac{-2+\sqrt{22}}{2}}$.
\end{loi_giai_cau}
